\documentclass[10pt,a4paper]{amsart}
\usepackage[margin=19mm]{geometry}
\usepackage{amsmath,amssymb,mathtools}
\usepackage[scr=rsfs,cal=euler]{mathalfa}
\usepackage{xfrac}
\usepackage[unicode,hidelinks,bookmarks=false]{hyperref}
\newcommand{\ie}{i.e.\ }
\newcommand{\cf}{cf.\ }
\newcommand{\resp}{respectively\ }
\newcommand{\Subsection}{\subsection}
\providecommand{\nobreakdash}{\nobreak\hbox{-}}
\setlength{\emergencystretch}{2em}
\multlinegap0pt
\usepackage{mathtools}
\usepackage{stmaryrd}
\usepackage{amssymb}
\usepackage{xcolor}
\usepackage[shortlabels]{enumitem}
\newtheorem{proposition}{Proposition}
\DeclareMathOperator{\gr}{gr}
\title{Note on a certain category of mod $p$ representations}
\author{Reinier Sorgdrager}
\date{Source: arXiv:2503.17773v2 (CC BY 4.0)}
\begin{document}
\maketitle

\noindent\emph{Source and licence.} Note on a certain category of mod $p$ representations. Version arXiv:2503.17773v2 (CC BY 4.0), distributed by arXiv under the Creative Commons Attribution 4.0 International licence, http://creativecommons.org/licenses/by/4.0/, as recorded in the arXiv licence metadata for this exact version and shown on the abstract page https://arxiv.org/abs/2503.17773v2. Full licence notice: \texttt{licenses/arxiv-2503.17773-CC-BY-4.0.txt}.\par\smallskip

\noindent\textit{Editorial scope.} Counted unit: the article's proposition centrality (source0.tex lines 157--159). The article's complete original proof of it is delivered with it (source0.tex lines 160--166, one \texttt{proof} environment of 7 lines). No mathematics is written, repaired or re-derived: the statement and the proof are the article's own lines, in the article's own notation, and no retained line is substituted or re-flowed.  Retained uncounted as the source-local context the proof needs: lines 122--124.  Nothing was omitted from any retained span, so the package carries no omission record.  No retained line contains a citation, so the wrapper carries no bibliography and no bibliography entry.  No retained line contains a figure environment, an image-inclusion command or drawing code, so no image is delivered and the package declares no asset dependency.  Editorial formatting only, in addition: an AMS \texttt{amsart} preamble that restates the article's own declarations and macros used by the retained lines, namely the environments \texttt{proposition} on separate counters, together with the standard packages those lines need.  The retained environments print in this document as Proposition 1 in order of appearance, whereas the article prints the same items with the numbering of its own full document.  The complete native source of the article as distributed in the arXiv e-print for this exact version is bundled unchanged as \texttt{sources/175205/source0.tex}.
	\begin{proposition}\label{maxideals}
		For $j\in\mathbf Z_{\geq0}$ we have $$\mathfrak m_G^j=\{x\in\mathbf F\llbracket G\rrbracket\colon\nu(x)\geq j\}.$$
	\end{proposition}

	\begin{proposition}\label{centrality}
		The ring $\gr_\nu\mathbf F\llbracket G^{p^N}\rrbracket=\mathbf F[a_0^{p^N},b_0^{p^N},c_0^{p^N},\dots,a_{f-1}^{p^N},b_{f-1}^{p^N},c_{f-1}^{p^N}]$ is a commutative polynomial $\mathbf F$-algebra. It is a subring of the center of $\gr\mathbf F\llbracket G\rrbracket$.
	\end{proposition}

	\begin{proof}
		The inclusion $\mathbf F\llbracket G^{p^N}\rrbracket\hookrightarrow\mathbf F\llbracket G\rrbracket$ induces an inclusion after taking $\gr_\nu$. We have $\gr_\nu\mathbf F\llbracket G\rrbracket=\gr\mathbf F\llbracket G\rrbracket$ by Proposition \ref{maxideals}.\\
		We first prove that the $a_i^{p^N}$ and $b_i^{p^N}$ are central. This follows from their centrality in $\gr\mathbf F\llbracket G\rrbracket$ which is seen as follows: take any generator $x\in\{a_0,b_0,\dots,a_{f-1},b_{f-1}\}$; then, as mentioned, $[a_i,x]$ lies in the span of the $c_j$; by the previous proposition $[a_i,x]$ is therefore central, from this and induction we get the formula $[a_i^\ell,x]=\ell a_i^{\ell-1}[a_i,x]$ for all $\ell\geq1$, since indeed we have \begin{equation*}
			[a_i^{\ell+1},x]=a_i[a_i^\ell,x]+a_ixa_i^\ell-xa_i^{\ell+1}=a_i[a_i^\ell,x]+[a_i,x]a_i^\ell
		\end{equation*} and to conclude from this last expression one uses the induction hypothesis and centrality of $[a_i,x]$; in particular, we see that $a_i^p$ is already central (the proof works the same for $b_i$ of course).\\
		Consider now the graded polynomial ring $\mathbf F[y_1,\dots,y_{3f}]$ which has the first $2f$ variables homogeneous in degree $1$ and the last $f$ variables homogeneous in degree $2$. It can be interpreted as the associated graded ring to $\mathbf F\llbracket p^N\mathbf Z_p^{3f}\rrbracket$ where $p^N\mathbf Z_p^{3f}$ is given the $p$-valuation induced from its homeomorphism to $G^{p^N}$. The homeomorphism yields an isomorphism of $\mathbf F$-modules $\mathbf F[y_1,\dots,y_{3f}]\xrightarrow{\sim}\gr_\nu\mathbf F\llbracket G^{p^N}\rrbracket$. By the previous proposition and the centrality of the $a_i^{p^N}$ and the $b_i^{p^N}$ this is actually a morphism of $\mathbf F$-algebras.
	\end{proof}

\end{document}
